\section{Conclusão}

\begin{frame}{Síntese dos Desafios e Soluções}
    \begin{itemize}
        \item As RSSF representam o pilar fundamental para a Internet das Coisas (IoT) e para a Indústria 4.0, superando as limitações físicas dos cabeamentos.  
        \item O maior gargalo técnico é a gestão energética: a sobrevivência de toda a malha depende da eficiência individual de cada sensor.  
        \item Nas camadas inferiores, protocolos de ciclo de trabalho (como o T-MAC) provaram ser indispensáveis para mitigar o desperdício crítico da escuta ociosa.  
        \item No roteamento, a transição para arquiteturas de múltiplos saltos (\textit{multi-hop}) e agregação de dados descentraliza a carga e estende a cobertura.
        \item A integridade do sistema é garantida por protocolos  
    \end{itemize}
\end{frame}

\begin{frame}{Considerações Finais}
    \begin{itemize}
        \item As redes evoluíram de sistemas dependentes de estações-base centralizadoras para organismos distribuídos.  
        \item O uso do aprendizado de máquina na borda (Edge AI) resolve o problema do desgaste da bateria evitando transmissões pesadas.  
    \end{itemize}

    \vspace{0.5cm}

    \begin{block}{O Caminho Definitivo}
        A integração entre a microeletrônica de baixíssimo consumo e a Inteligência Artificial de borda é o caminho consolidado para sistemas de monitoramento autônomos e economicamente viáveis.  
    \end{block}
    
    \vspace{0.5cm}
    

\end{frame}